\section*{Phần B --- Quy trình làm bài theo dạng (bám các lời giải .tex trong repo)}

\subsection*{B.1. Chương 2 --- Tiền xử lý}
\begin{enumerate}
  \item Sắp các giá trị một chiều tăng dần.
  \item Tính trung bình, trung vị, mốt; ghi rõ có nhiều mốt nếu các giá trị có cùng tần suất cực đại.
  \item Tính \(Q_1\), \(Q_3\), \(\mathrm{IQR}\); ghi rõ quy ước vị phân vị (vị trí nguyên hoặc nội suy giữa hai quan sát kề) vì các bài ôn có thể dùng hai cách.
  \item Boxplot \& ngoại lai Tukey: năm số tóm tắt, \(\mathrm{IQR}\), so sánh với \(Q_{1,3}+1{,}5\,\mathrm{IQR}\).
  \item \emph{Rời rạc hoá chia đều tần số} (\(k\) nhóm): lấy dãy đã sắp, chia thành các khúc có số điểm gần bằng nhau như các bài ôn (ví dụ \(q{+}\cdots\) với chia nguyên của \(n/k\)).
    \footnote{Các mục công thức chính xem Phần~A.}
  \item Thang đo khoảng: chứng minh tồn tại \(a{>}0\), \(b\) sao cho \(W=aV+b\) trên các cặp được cho hoặc bác bỏ nếu dãy không cùng phép affine tăng.
\end{enumerate}

\subsection*{B.2. Chương 3 --- Đánh giá mô hình}
\begin{enumerate}
  \item Nếu đề cho dãy luật có thứ tự \texttt{if}-\texttt{elseif}: áp đúng thứ tự trên mỗi mẫu rồi đối chiếu nhãn thực.
  \item Lập confusion matrix (hàng: thực, cột: đoán) bằng cách đếm từng cặp \((y_i,\hat{y}_i)\).
  \item Nhị phân với một lớp dương quy định: đếm \(\mathrm{TP},\mathrm{FN},\mathrm{FP},\mathrm{TN}\) rồi thế vào các công thức ở Phần~A.
  \item Đa lớp: accuracy tổng thể \(\sum_k \mathrm{TP}_k / n\) và chỉ số theo lớp \(k\) như trong Phần~A.\@
\end{enumerate}

\subsection*{B.3. §\,4.1 --- Cây quyết định}
\begin{enumerate}
  \item Tại nút \(S\): ghi các mẫu và các đặc trưng được phép chia.
  \item Với từng đặc trưng có giá trị rời rạc trong \(S\), tách thành vài nhánh \(\{S_v\}\) theo đủ các giá trị xuất hiện.
  \item Tính chỉ báo của đề: Entropy của \(S\), Gain cho mỗi đặc trưng \(\mathrm{Gain}(S,A)=H(S)-\sum_v \frac{|S_v|}{|S|}\,H(S_v)\); hoặc Gini có trọng số; hoặc lỗi phân có trọng số.
  \item Chọn thuộc tính theo chỉ báo đề cho: \(\mathrm{Gain}\) lớn nhất hoặc Gini có trọng số nhỏ nhất hoặc lỗi có trọng số nhỏ nhất.
  \item Lặp lại trên từng nhánh chứa mẫu chưa thuần nhãn \(\land\) vẫn còn đặc trưng dùng được; đặt lá khi chỉ một nhãn còn lại (hoặc hết chia theo ôn tay).
\end{enumerate}
\noindent\textbf{Dự báo mẫu mới:} đưa các giá trị đặc trưng vào các luật if--then từng bước từ gốc của cây xuống lá. Trường hợp đề chỉ cho dãy luật if--elseif có thứ tự cố định thì luôn so khớp từ trên xuống và dừng ở luật đầu thoả như phần Đánh giá.

\subsection*{B.4. §\,4.2 --- Hồi quy tuyến tính + batch GD}
\begin{enumerate}
  \item Xác định \(m\) và \(\hat{y}_i=wx_i+b\); \(J=\dfrac{1}{2m}\sum_i (y_i-\hat{y}_i)^2\).
  \item Tại \((w^{(t)}, b^{(t)})\): tính các dự báo \& phần dư \(y_i-\hat{y}_i\).
  \item \(\partial J/\partial w\), \(\partial J/\partial b\) theo Phần~A.
  \item Cập nhật \(w\leftarrow w-\eta\,\partial J/\partial w\), \(b\leftarrow b-\eta\,\partial J/\partial b\) (\(\eta\) của đề; đổi \(\eta\) \(\leftrightarrow\) \(\rho\) theo notation đề ôn tay).
\end{enumerate}

\subsection*{B.5. §\,4.3 --- Phân lớp tuyến tính (sign) + hinge trên tập sai + Perceptron}
\begin{enumerate}
  \item Tính \(s_k=w^\top x_k+b\) rồi \(\hat{y}_k=\mathrm{sign}(s_k)\); so với \(t_k\) để tìm các mẫu sai.
  \item Chi phí chỉ trên tập sai \(\mathcal{Y}\): lấy \(\partial J/\partial w\), \(\partial J/\partial b\) theo Phần~A.
  \item Với gradient trên tập sai (§\,4.6.3): \(w\leftarrow w-\rho\,\partial J/\partial w\), \(b\leftarrow b-\rho\,\partial J/\partial b\).
  \item Với Perceptron một mẫu (§\,4.3): chọn mẫu sai theo thứ tự đề; \(w\leftarrow w+\rho\,t\,x\), \(b\leftarrow b+\rho\,t\).
\end{enumerate}

\subsection*{B.6. §\,4.4 --- \(k\)-NN}
\begin{enumerate}
  \item Mã hoá thuộc tính rời rạc theo quy ước đề (vd. Có/Không \(\to 0/1\)); xây \(\mathbf{x},\mathbf{q}\) cùng số chiều.
  \item Tính \(d^2(\mathbf{x},\mathbf{q})\) rồi sắp tăng \(d\); lấy \(K\) láng giềng nhỏ nhất.
  \item Bỏ phiếu nhãn đa số; có thể nhắc lại với \(K'\) khác nếu đề hoặc nhận xét yêu cầu.
\end{enumerate}

\subsection*{B.7. §\,4.5 --- Naive Bayes + Laplace}
\begin{enumerate}
  \item Đếm \(n,n_c\) và số nhãn; Laplace \(\alpha{=}1\): \(P(c)=\dfrac{n_c+\alpha}{n+\alpha|\mathcal{C}|}\); với thuộc tính rời rạc có \(V\) giá trị khả dụng, dùng \(\dfrac{n_{v,c}+1}{n_c+V}\) với \(n_{v,c}\) là số lần thấy \(v\) trong lớp \(c\) (theo các ôn 4.5).
  \item Gaussian: tính \(\mu_c,s_c\) trong từng lớp; lấy log mật độ tại mỗi biến liên tục.
  \item Nhận các thể hiện mới: cộng \(\log P(c)+\sum \log P(v_j\mid c)\) (bỏ qua thuộc tính thiếu nếu đề nói không quan sát).
  \item Chọn lớp có log-score lớn hơn.
  \item Multinomial (email): tiên nghiệm lớp; mỗi token \(P(w_i\mid c)=(\text{đếm}+1)/(N_c+V)\); nhân (log) theo tần suất bộ đếm của văn bản cần phân lớp.
\end{enumerate}

\subsection*{B.8. §\,4.6 --- Perceptron / MLP tí hon}
\begin{enumerate}
  \item Viết đường biên dạng \(\mathbf{w}^\top \mathbf{x}+b=0\) từ đề đồ họa/phương trình.
  \item Thay điểm vào \(\tau(s)\): \(+1\) nếu \(s\ge 0\) ngược lại \(-1\) (hay quy định đề).
  \item MLP không ẩn: tính tuần tự theo chỉ báo các nút từ trái-phải theo các trọng số đề.
  \item OR/AND trong \(\{0,1\}^2\): chỉ ra \(\mathbf{w},b\) sao \(\tau(\mathbf{w}^\top\mathbf{x}+b)\) khớp bảng (so với ví dụ \(-\tfrac12+x_1+x_2\) cho OR trong lời giải đã ôn).
\end{enumerate}

\subsection*{B.9. Chương 5 --- \(K\)-means}
\begin{enumerate}
  \item Khởi tạo các tâm \(C_j^{(0)}\) theo đề (\(j{=}1,\ldots,K\)).
  \item Bước \emph{gán}: mỗi điểm vào cụm của tâm gần Euclidean nhất (nếu hai khoảng cách bằng nhau, ghi rõ quy tắc phá hoà, vd. chọn cụm có chỉ số nhỏ hơn).
  \item Bước \emph{cập nhật}: tâm mới = trung bình toạ độ của điểm trong cụm.
  \item Lặp đến khi gán không đổi hoặc các tâm không đổi (trên bài tay); ghi kết quả cụm cuối \& tâm.
\end{enumerate}

\subsection*{B.10. Chương 5 --- Phân cụm phân cấp (single linkage)}
\begin{enumerate}
  \item Tính ma trận khoảng cách đôi một giữa các điểm.
  \item Tìm cặp có khoảng cách nhỏ nhất để ghép thành cụm ở bước hiện tại; độ cao nối trên dendrogram thường bằng khoảng cách của bước đó trong các bài ôn.
  \item Sau khi ghép hai cụm \(U\) và \(V\), với mọi cụm \(W\) còn lại: cập nhật
  \[
    d_{\mathrm{single}}(U{\cup}V,W)=\min\bigl\{d_{\mathrm{single}}(U,W),\, d_{\mathrm{single}}(V,W)\bigr\}.
  \]
  \item Lặp cho đến khi còn một cụm.
\end{enumerate}


